% =============================================================================
% bab3-metode-bisnis.tex — BAB III METODE ANALISIS BISNIS
% Proyek Bisnis / Studi Kasus — Lampiran II (S2)
% =============================================================================

\chapter{METODE ANALISIS BISNIS}
\label{bab:metode-bisnis}

\section{Jenis Penelitian}
\label{sec:jenis-penelitian-bisnis}

Penelitian ini merupakan [proyek bisnis/studi kasus] dengan pendekatan
[kualitatif/kuantitatif/campuran]. Uraikan jenis penelitian, paradigma yang digunakan,
serta alasan pemilihan pendekatan tersebut.

Penelitian ini bersifat [deskriptif/eksploratoris/eksplanatoris] yang bertujuan untuk
[mendeskripsikan/mengeksplorasi/menjelaskan] kondisi [aspek bisnis] pada [nama organisasi].

\section{Jenis dan Sumber Data}
\label{sec:jenis-sumber-data}

\subsection{Jenis Data}
Data yang digunakan dalam penelitian ini terdiri dari:
\begin{enumerate}
  \item \textbf{Data Primer:} Data yang diperoleh langsung dari narasumber/objek
        penelitian melalui wawancara mendalam, observasi, dan/atau kuesioner.
  \item \textbf{Data Sekunder:} Data yang diperoleh dari laporan keuangan, laporan
        tahunan, dokumen internal perusahaan, publikasi ilmiah, dan sumber-sumber
        sekunder lainnya.
\end{enumerate}

\subsection{Sumber Data}
Sumber data dalam penelitian ini adalah:
\begin{enumerate}
  \item Internal: Manajemen dan karyawan [nama organisasi], dokumen internal perusahaan.
  \item Eksternal: Publikasi industri, laporan pasar, data BPS, dan penelitian terdahulu.
\end{enumerate}

\subsection{Teknik Pengumpulan Data}
Pengumpulan data dilakukan melalui:
\begin{enumerate}
  \item \textbf{Wawancara mendalam (in-depth interview):} dilakukan kepada \ldots
  \item \textbf{Observasi:} pengamatan langsung terhadap \ldots
  \item \textbf{Studi Dokumentasi:} analisis dokumen-dokumen \ldots
  \item \textbf{Kuesioner} (jika ada): diberikan kepada \ldots
\end{enumerate}

\section{Teknik Analisis Data}
\label{sec:teknik-analisis-bisnis}

\subsection{Alat Analisis}
Alat analisis yang digunakan dalam penelitian ini adalah [nama alat analisis],
antara lain:

\begin{enumerate}
  \item \textbf{Analisis SWOT:} Digunakan untuk mengidentifikasi kekuatan,
        kelemahan, peluang, dan ancaman organisasi.
  \item \textbf{[Alat Analisis Lain]:} [Penjelasan singkat].
\end{enumerate}

\subsection{Tahapan Analisis}
Proses analisis data dilakukan melalui tahapan berikut:
\begin{enumerate}
  \item Pengumpulan dan triangulasi data.
  \item Reduksi data dan identifikasi tema.
  \item Analisis menggunakan alat analisis yang telah ditentukan.
  \item Perumusan rekomendasi solusi bisnis.
  \item Validasi temuan dan rekomendasi.
\end{enumerate}
